\section{One coefficient request for all the marked cover levels}

% BEGIN U76W1
\paragraph{U76W1. A finite set of algebraic numbers at a prescribed root and level.}
Fix a complex number $r$ and an integer $m>0$. There is a finite
set $Z(r,m)$ of algebraic numbers such that for every permitted CM
host with marked root $r$, every intermediate field of $\mathbb C/\mathbb Q$
containing that set satisfies the source condition $\mathrm{WeberIn}$ at $m$.

If there is no such host, take the empty set. Otherwise the periods of
every such host form exactly the same lattice
$L=\mathbb Z+\mathbb Z(1+r)/2$. Write $g_2,g_3$ for its invariants
and $\Delta=g_2^3-27g_3^2$. Take
\[
 Z(r,m)=\{g_2^3/\Delta\}\ \cup \left\{\frac{g_2g_3}{\Delta}\wp_L(z):
               0\ne z\in m^{-1}L/L\right\}.
\]
Here $\wp_L$ is periodic, so choosing representatives does not affect
the set. The source quantifies over period pairs, and their ordered
bases may differ. Their Weierstrass functions and invariants depend
only on the lattice. The $m$-torsion has $m^2$ elements, giving
at most $m^2$ requested numbers, including when $m=1$.

For algebraicity, multiplication by $(1+r)/2$ preserves $L$ and is
not an integer. The singular-modulus theorem for a lattice with
complex multiplication gives $j(L)=1728g_2^3/\Delta$ algebraic.
Choose $u$ with $u^{12}=\Delta\ne0$. On the scaled lattice $uL$
the discriminant is one, the cube of its $g_2$ is $j(L)/1728$,
and the square of its $g_3$ is $(j(L)/1728-1)/27$. The scaled
Weierstrass cubic is thus defined over $\overline{\mathbb Q}$.
Every point of its finite $m$-torsion is algebraic over that field.
At a nonzero torsion point its $x$-coordinate is $\wp_{uL}(uz)$.
Multiplying by the two scaled invariants gives precisely the
original normalized Weber value, by the homothety laws of weights
$4,6,2$. Thus all requested numbers are algebraic. Containment
of this set is exactly the two required inclusions for every marked
period pair.
% END U76W1

% BEGIN U76W2
\paragraph{U76W2. The field request precedes the choice of a tuple.}
For a prehost $P$ and a prescribed character family $\Xi$, let
$A_\Xi$ be the set of finite primes at which at least one of these
characters ramifies. When this set is finite, take all its
subsets $S$ and put
\[
 m_S=g_0g_1\prod_{q\in\{\#k_Q:Q\in S\}}q,\qquad
 Z=\bigcup_{S\subseteq A_\Xi}
 Z\bigl(\phi_0(\sqrt{-d}),m_S\bigr).
\]
Each distinct residue-field cardinality occurs once. Each level is
positive, and U76W1 makes $Z$ a finite set of algebraic numbers.
For every $S$ whose primes are witnessed by characters in $\Xi$,
the inclusion $S\subseteq A_\Xi$ puts its entire Weber set in $Z$.
Consequently any coefficient field containing $Z$ works at every
one of these levels and for every host with the prescribed root.

This gives the original eventual field condition
$\mathrm{large}_2(K,\mathrm{WeberLev})$ for every context $K$.
Use $Z$ for its fourth, complex-algebraic request and the empty
sets for its first three requests. On arbitrary test families with
infinite $A_\Xi$, prescribe the empty fourth set: its specification
there is an implication with finite $A_\Xi$ as antecedent.
On an actual admissible family $\Xi=T.\mathrm{set}(f)$, the finite
twist-family law and the finite ramification of each character make
$A_\Xi$ finite. The field-containment clause of the size condition
therefore gives the assertion just proved for its own embedded
coefficient field $B.\mathrm{LLC}$. It is one request per prehost
and character family, before the eventual column or its tuple is
specified.
% END U76W2

% BEGIN U76W3
\paragraph{U76W3. The cover-field input on the selected column.}
Take a base on which U76W2's $\mathrm{WeberLev}$ holds, and fix an
original column $w$, its level and any three-stage geometric record
$z$ of that column. Then $\mathrm{CoverField}(z.\mathrm{core},R)$
holds, where $R=z.R.\mathrm{toRealDat}$. To prove it take any
row $i$, tuple $t$ and higher model in this condition. Each
$Q\in z.R.Q(t)$ is a ramified prime of the character $\xi_t^H$;
the specified range of the column puts that character in
$T.\mathrm{set}(f)$. Apply the single field condition to
$S=z.R.Q(t)$. The geometric pins identify its marked root
with $\phi_0(\sqrt{-d})$, its lower level with $g_0g_1$, and the
field required for this higher model with $B.\mathrm{LLC}$.
Its level is exactly $m_S$: the extra factor in the level is
the product of the distinct residue cardinalities used above.
Its auxiliary level parameter is one, so the effective level
$\operatorname{lcm}(m_S,1)$ equals $m_S$. The model has level
at least three. Apply the marked-component argument of W76F3 on
this same model at $m_S$. It gives the required relative
algebraic closure in its component function fields. These
choices were arbitrary in $\mathrm{CoverField}$.

In particular, suppose the given $z$ also has its uniform
rational coordinates $\mathrm{LUnif}$, oscillator extension
$\mathrm{IoExt}$ and its base-row maps $\mathrm{OneMapRowR}$
with constant phase one. Its geometric record supplies the
weight and frame tie, conjugation, tuple pins and phase law.
The phase prescribed by the three-stage test is one as well:
$z.\mathrm{pins}$ gives $z.R.\mathrm{phs}_i=1$ and every
field automorphism fixes one. These fields and the cover-field
conclusion give the full $\mathrm{TiesG}$ of T76R1. With the
corresponding $\mathrm{TiesB}$, T76R4 now gives the rational
lift tests and the integral cusp lattices on this column.
The original selection problem still calls for the geometric
record and those rational coordinate and base-map inputs on the
same eventual family.
% END U76W3
